% Treść — 2 kolumny (multicol).
%
% Dostępne makra brandowe:
%   \glitchline{tekst}  — chromatic aberration na fragmencie tekstu
%   \corr               — inline blok "corrupted" (czerwony marker)
%   \corrk              — inline blok "corrupted" (ciemny marker)
%   \slaycross          — mini celownik
%   \disruptline        — linia szumu między sekcjami
%   \noisebar[N]{label} — pasek statusu N% z labelką
%   \slaybarcode[kod]   — pseudo-barcode z etykietą
%
% Sekcje mają automatyczny hex prefix (0x01, 0x02...) + glitch.
% Użyj \subsection{} dla podsekcji z mniejszym hex prefixem.

\begin{multicols}{2}

\section{Wprowadzenie}

Po co ten dokument i jakie pytanie stawiamy. Jedno zdanie tezy +
jak ją zweryfikujemy (pomiar). Powiązanie z nitką dialektyczną
(Teza $\to$ Synteza).

\section{Metoda}

Skąd dane, jak je przetwarzamy, jaki model/eksperyment.
Definicje formalnie:
\[
  \mathrm{metryka} = \frac{1}{N}\sum_i f(x_i)
\]

\section{Wyniki}

Tu ląduje empiria. Przykładowa tabela:

\begin{table}[h]
  \centering
  \caption{Przykładowe wyniki.}
  \begin{tabular}{lrr}
    \toprule
    Wariant & metryka $\downarrow$ & próbek \\
    \midrule
    baseline   & 1.80 & --- \\
    nasz model & \glitchline{1.20} & --- \\
    \bottomrule
  \end{tabular}
\end{table}

\disruptline

\section{Warunek obalenia}

% Bramka empiryczna — NIE usuwaj tej sekcji.
% Każda synteza musi nieść test/liczbę/próg, który mógłby ją wywrócić.
Co konkretnie sfalsyfikowałoby tezę: test, liczba albo próg.
Np. „jeśli metryka na held-oucie $>X$, teza upada".\,\corr

\section{Dyskusja}

Co wynik oznacza, gdzie się nie broni, co dalej.
Uczciwie o ograniczeniach.

\vspace{6pt}
\noisebar[100]{DOKUMENT·KOMPLETNY}

\vspace{4pt}
{\fontsize{5.5}{5.5}\selectfont\ttfamily\color{BETON}%
  \glitchline{SL/AI·LAB}~·~WROCŁAW~·~OFFLINE~BY~CHOICE~\slaycross[BETON!60]}

\end{multicols}
